% 一维高斯波包的时间—能量不确定度的计算
% license Usr
% type Note

\subsection{粒子}
\begin{equation}
P(k)\dd{k} = P'(E)\dd{E} = P'(k^2/2)k\dd{k}~.
\end{equation}
所以
\begin{equation}
P'(E) = P(k)/k = \exp[-a^2 k^2]/k~.
\end{equation}

\begin{equation}
\Delta k = \Delta E/k~.
\end{equation}
\begin{equation}
\Delta x = c\Delta t~.
\end{equation}


高斯波包 FWHMI （用 $\Delta$ 表示）和 sigma 的关系是？前者是后者的 $2\sqrt{2\ln 2} = 2.35482$ 倍。

$\frac{c\Delta t \cdot \Delta E/k}{8\ln 2} = \frac{\Delta x}{2\sqrt{2\ln 2}}\frac{\Delta k}{2\sqrt{2\ln 2}} = \sigma_x \sigma_k \leq \frac{1}{2}$

于是得到（原子单位） $\Delta t\Delta E \leq (4\ln 2)k/c = 2.77258 k/c$。

\begin{lstlisting}[language=matlab]
Eng = eV2au_E(143); % 83
FWHMI_E = eV2au_E(20); % 10

k = sqrt(2*Eng);
c = 137.036;

FWHMI_t = 4*log(2)*k/c / FWHMI_E;

disp('FWHMI_t [as] = ');
disp(au_t2as(FWHMI_t));
\end{lstlisting}

\subsection{电磁波}
$E = \omega$。

\begin{equation}
\frac{\Delta t}{2\sqrt{2\ln 2}}\frac{\Delta E}{2\sqrt{2\ln 2}} = \sigma_t \sigma_\omega \leq \frac{1}{2}~.
\end{equation}
\begin{equation}
\Delta t\Delta E = 4\ln 2 = 2.772589~.
\end{equation}

国际单位下，$\beta_E\beta_t = \hbar$，所以最终换算成电子伏和阿秒为单位，有
\begin{equation}
\Delta t \Delta E  = \frac{2\ln 2}{Q_E\pi}h = 1.824934\e{3}~.
\end{equation}
